% ============================================================================
% LANGENTWURF LE-015: Messenger
% Profil: S (Senioren 65+) | Tier: 2 (Standard)
% ============================================================================

\documentclass[11pt,a4paper]{article}
\usepackage[utf8]{inputenc}
\usepackage[T1]{fontenc}
\usepackage[ngerman]{babel}
\usepackage{geometry}
\usepackage{fancyhdr}
\usepackage{lastpage}
\usepackage{xcolor}
\usepackage{tcolorbox}
\usepackage{enumitem}
\usepackage{longtable}
\usepackage{hyperref}
\usepackage{amssymb}

\geometry{left=2.5cm, right=2.5cm, top=2.5cm, bottom=2.5cm}
\definecolor{fach}{HTML}{b35c00}
\definecolor{gray}{HTML}{666666}
\definecolor{successgreen}{HTML}{2a7a4b}
\definecolor{warnred}{HTML}{c62828}

\tcbuselibrary{skins,breakable}
\newtcolorbox{scorebox}{ colback=white, colframe=fach, fonttitle=\bfseries, title=Didaktik-Score, breakable }
\newtcolorbox{warnbox}[1][]{ colback=warnred!5, colframe=warnred, fonttitle=\bfseries, title=#1, breakable }

\pagestyle{fancy}
\fancyhf{}
\fancyhead[L]{\small\textcolor{gray}{Digitale Grundlagen -- 015 Messenger}}
\fancyhead[R]{\small\textcolor{gray}{LE Tier 2 / Profil S}}
\fancyfoot[C]{\small\thepage{} / \pageref{LastPage}}
\renewcommand{\headrulewidth}{0.4pt}

\begin{document}

\begin{titlepage}
    \centering
    \vspace*{2cm}
    {\Large\textcolor{gray}{Digitale Grundlagen -- Senioren 65+}}\\[2cm]
    {\Huge\textcolor{fach}{\textbf{Langentwurf}}}\\[0.5cm]
    {\large Tier 2 | Profil S}\\[2cm]
    {\LARGE\textbf{015 -- Messenger: WhatsApp \& Co.}}\\[0.5cm]
    {\large Modul C: Internet \& Kommunikation}\\[2cm]
    \begin{tabular}{rl}
        \textbf{Zeitrahmen:} & 75--90 Minuten \\
        \textbf{Vorkenntnisse:} & E-Mail, Kontakte anlegen \\
    \end{tabular}
    \vfill
    {\normalsize Stand: \today}
\end{titlepage}

\tableofcontents
\newpage

\section{Bedingungsanalyse}

\subsection{Ausgangssituation}
\begin{itemize}
    \item WhatsApp ist oft der Hauptgrund für Smartphone-Kauf
    \item Familie und Freunde nutzen es -- ``Ich muss erreichbar sein''
    \item Unterschied zur SMS unklar
    \item Gruppen sind überfordernd (zu viele Nachrichten)
    \item Sicherheitsbedenken: ``Liest WhatsApp alles mit?''
\end{itemize}

\subsection{Typische Probleme}
\begin{itemize}
    \item ``Wie finde ich meine Kontakte in WhatsApp?''
    \item ``Die Gruppe schreibt so viel, ich finde nichts mehr''
    \item ``Wie schicke ich ein Foto?''
    \item ``Jemand hat mir geschrieben -- aber wer ist das?''
\end{itemize}

\section{Sachanalyse}

\subsection{Kernkonzepte}
\begin{enumerate}
    \item \textbf{Messenger:} App für schnelle Nachrichten über Internet (nicht SMS)
    \item \textbf{Chat:} Gespräch mit einer Person oder Gruppe
    \item \textbf{Gruppe:} Mehrere Personen in einem Chat zusammen
    \item \textbf{Status:} Bilder/Texte, die nach 24 Stunden verschwinden
    \item \textbf{Haken:} Ein Haken = gesendet, zwei Haken = angekommen, blaue Haken = gelesen
\end{enumerate}

\subsection{Alltagsanalogien}
\begin{tabular}{|l|p{8cm}|}
\hline
\textbf{Begriff} & \textbf{Analogie} \\
\hline
Messenger & Digitales Gespräch -- wie telefonieren, nur schriftlich \\
Chat & Ein Gespräch mit einer Person \\
Gruppe & Runder Tisch, an dem alle mithören und mitreden \\
Sprachnachricht & Kurzer Anrufbeantworter, den du jederzeit abhören kannst \\
Status & Aushang am schwarzen Brett (für alle, zeitlich begrenzt) \\
\hline
\end{tabular}

\section{Didaktische Analyse}

\subsection{Stellung in der Reihe}
Alternative/Ergänzung zu E-Mail. Wichtigste alltägliche Kommunikation für viele Senioren.

\subsection{Didaktische Reduktion}
\textbf{Ausgeklammert:} Anrufe über WhatsApp, Datenschutz-Einstellungen im Detail, andere Messenger.

\textbf{Fokus:} WhatsApp als Hauptbeispiel, Nachrichten und Fotos senden/empfangen.

\subsection{Sicherheitshinweis}
\begin{warnbox}[Vorsicht bei unbekannten Nachrichten!]
\begin{itemize}
    \item ``Hallo Oma, ich habe eine neue Nummer'' -- IMMER nachfragen!
    \item Unbekannte Nummern? Nicht auf Links klicken
    \item Im Zweifel: Anrufen und nachfragen (alte bekannte Nummer)
\end{itemize}
\end{warnbox}

\section{Lernziele}

\begin{tabular}{|c|p{7cm}|c|c|}
\hline
\textbf{Nr.} & \textbf{Die Teilnehmer können...} & \textbf{Stufe} & \textbf{Niveau} \\
\hline
FZ1 & WhatsApp öffnen und die Übersicht verstehen & Anwenden & Basis \\
FZ2 & eine Nachricht an einen Kontakt schreiben & Anwenden & Basis \\
FZ3 & ein Foto aus der Galerie senden & Anwenden & Basis \\
FZ4 & die Haken-Bedeutung erklären & Verstehen & Basis \\
FZ5 & verdächtige Nachrichten erkennen (Enkel-Trick) & Anwenden & Vertiefung \\
\hline
\end{tabular}

\section{Verlaufsplanung}

\begin{longtable}{|p{1cm}|p{2cm}|p{5cm}|p{3cm}|p{2cm}|}
\hline
\textbf{Zeit} & \textbf{Phase} & \textbf{Inhalt} & \textbf{TN-Aktivität} & \textbf{Material} \\
\hline
0--10 & Einstieg & ``Wer nutzt WhatsApp?'' Erfahrungen sammeln & Berichten & -- \\
\hline
10--25 & Übersicht & App öffnen, Chats-Liste, einzelner Chat, Eingabefeld & Mitmachen & S.1 \\
\hline
25--40 & Nachricht & Text tippen, Emojis, Senden, Haken verstehen & Üben & S.1 \\
\hline
40--55 & Fotos & Foto aus Galerie senden, Kamera-Symbol & Selbst versuchen & S.2 \\
\hline
55--75 & Gruppen & Gruppen verstehen, Benachrichtigungen stummschalten & Eigene Gruppen & S.2 \\
\hline
75--85 & Sicherheit & Enkel-Trick digital, verdächtige Nachrichten & Beispiele & S.3 \\
\hline
\end{longtable}

\section{Lösungen}

\textbf{WhatsApp-Bereiche:}
\begin{itemize}
    \item \textbf{Chats:} Liste aller Gespräche (Personen und Gruppen)
    \item \textbf{Status:} Bilder/Videos, die nach 24h verschwinden (Facebook-Stories)
    \item \textbf{Anrufe:} Telefonate über WhatsApp (braucht WLAN/mobile Daten)
\end{itemize}

\textbf{Nachricht senden:}
\begin{itemize}
    \item Chat öffnen (Person antippen)
    \item Unten ins Textfeld tippen
    \item Text eingeben
    \item Grüner Pfeil = Senden
\end{itemize}

\textbf{Foto senden:}
\begin{itemize}
    \item Im Chat: + Symbol (links) oder Kamera (rechts)
    \item ``Foto- und Videomediathek'' wählen
    \item Foto auswählen, optional Text hinzufügen
    \item Senden
\end{itemize}

\textbf{Die Haken:}
\begin{itemize}
    \item $\checkmark$ Ein grauer Haken: Nachricht wurde gesendet
    \item $\checkmark\checkmark$ Zwei graue Haken: Nachricht ist angekommen
    \item \textcolor{blue}{$\checkmark\checkmark$} Zwei blaue Haken: Nachricht wurde gelesen
\end{itemize}

\section{Didaktik-Score}

\begin{scorebox}
\begin{tabular}{|c|l|c|c|p{3.5cm}|}
\hline
\# & Dimension & Gew. & Score & Notiz \\
\hline
1 & Differenzierung & 15\% & 8/10 & Basis/Vertiefung gut \\
2 & Taxonomie & 10\% & 9/10 & Verstehen + Anwenden \\
3 & Alltagsrelevanz & 10\% & 10/10 & Hauptkommunikation! \\
4 & Aktivierung & 10\% & 10/10 & Selbst schreiben \\
5 & Scaffolding & 10\% & 9/10 & Schritt-für-Schritt \\
6 & Feedback & 10\% & 8/10 & Haken erklären \\
7 & Sprachliche Klarheit & 10\% & 9/10 & Gute Analogien \\
8 & Motivation & 10\% & 10/10 & Familie erreichen \\
9 & Barrierefreiheit & 10\% & 9/10 & Große Darstellung \\
10 & Praktikabilität & 5\% & 8/10 & 85 Min evtl. lang \\
\hline
\multicolumn{3}{|r|}{\textbf{GESAMT:}} & \textbf{9.15/10} & \textcolor{successgreen}{\textbf{FREIGABE}} \\
\hline
\end{tabular}
\end{scorebox}

\end{document}
